\section{Introduction et objectifs}

\subsection{Contexte}
Dans un climat chaud comme celui de Dakar, la chaleur reçue par l'enveloppe d'un bâtiment (air extérieur
chaud, soleil sur le toit) se traduit par une surchauffe intérieure et par des besoins de refroidissement.
Choisir l'isolation, la couleur du toit ou la stratégie de ventilation demande de comparer des solutions
de façon quantitative, avec un modèle dont on sait qu'il est juste.

\subsection{Objectif général}
Développer un outil numérique de simulation thermique fondé sur la modélisation mathématique et la
méthode des éléments finis, puis l'utiliser pour comparer des configurations de bâtiment.

\subsection{Objectifs spécifiques}
\begin{itemize}
  \item définir une géométrie 2D simplifiée et les matériaux ;
  \item établir les équations de transfert thermique et les conditions aux limites ;
  \item construire un maillage adapté et résoudre le problème par éléments finis P1 avec FreeFem++ ;
  \item exporter les résultats vers Python pour les traiter et les visualiser ;
  \item comparer plusieurs scénarios et produire un rapport synthétique ;
  \item développer une interface permettant de modifier les paramètres.
\end{itemize}

\subsection{Démarche}
Le projet avance par petites étapes, chacune validée avant la suivante : un modèle physique simple est
d'abord vérifié contre des solutions connues, puis enrichi (plusieurs matériaux, temps, pièce en 2D,
air intérieur, soleil et ventilation). Une image colorée ne prouve pas que le calcul est juste : chaque
cas est accompagné de vérifications quantitatives (solution exacte, étude de convergence, bilans de flux).

\subsection{Limites annoncées}
Le modèle est 2D et simplifié. Les résultats dépendent de la qualité des données et des hypothèses.
Ils constituent une aide à l'étude, non une certification réglementaire. Les limites sont détaillées
à la section~\ref{sec:limites}.
